% from ega3-1-fr.tex
\oldpage[III]{119}
\subsection{Le théorème de finitude : cas des schémas formels}
\label{subsection:III.3.4-fr}

Les résultats de cette section (sauf la définition
(\hyperref[III.3.4.1-fr]{3.4.1})) ne seront pas utilisés dans la suite de ce
chapitre.

\begin{env}[3.4.1]
\label{III.3.4.1-fr}
Soient $\mathfrak X$ et $\mathfrak S$ deux préschémas formels localement
noethériens (I, 10.4.2), $f:\mathfrak X\to\mathfrak S$ un morphisme de
préschémas formels. Nous dirons que $f$ est un morphisme \emph{propre} s'il
vérifie les conditions suivantes~:
\begin{enumerate}
\item[$1^\circ$] $f$ est un morphisme de type fini (I, 10.13.3).
\item[$2^\circ$] Si $\mathcal K$ est un Idéal de définition de $\mathfrak S$
et si l'on pose
$\mathcal J=f^*(\mathcal K)\mathcal O_{\mathfrak X}$,
$X_0=(\mathfrak X,\mathcal O_{\mathfrak X}/\mathcal J)$,
$S_0=(\mathfrak S,\mathcal O_{\mathfrak S}/\mathcal K)$, le morphisme
$f_0:X_0\to Y_0$ déduit de $f$ (I, 10.5.6) est propre.
\end{enumerate}

Il est immédiat que cette définition ne dépend pas de l'Idéal de définition
$\mathcal K$ de $\mathfrak S$ considéré~; en effet, si $\mathcal K'$ est un
second Idéal de définition tel que $\mathcal K'\subset\mathcal K$, et si on
pose $\mathcal J'=f^*(\mathcal K')\mathcal O_{\mathfrak X}$,
$X'_0=(\mathfrak X,\mathcal O_{\mathfrak X}/\mathcal J')$,
$S'_0=(\mathfrak S,\mathcal O_{\mathfrak S}/\mathcal K')$, le morphisme
$f'_0:X'_0\to S'_0$ déduit de $f$ est tel que le diagramme
\[
\xymatrix{
X_0\ar[r]^{f_0}\ar[d]_i & S_0\ar[d]^j\\
X'_0\ar[r]_{f'_0} & S'_0
}
\]
est commutatif, $i$ et $j$ étant des immersions surjectives~; il revient donc
au même de dire que $f_0$ ou $f'_0$ est propre, en vertu de (II, 5.4.5).

On notera que, pour tout $n\geq0$, si on pose
$X_n=(\mathfrak X,\mathcal O_{\mathfrak X}/\mathcal J^{n+1})$,
$S_n=(\mathfrak S,\mathcal O_{\mathfrak S}/\mathcal K^{n+1})$, le morphisme
$f_n:X_n\to Y_n$ déduit de $f$ (I, 10.5.6) est propre pour tout $n$ dès qu'il
l'est pour $n=0$ (II, 5.4.6).

Si $g:Y\to Z$ est un morphisme propre de préschémas usuels, localement
noethériens, $Z'$ une partie fermée de $Z$, $Y'$ une partie fermée de $Y$
telle que $g(Y')\subset Z'$, le prolongement
$\widehat g:Y_{/Y'}\to Z_{/Z'}$ de $g$ aux complétés (I, 10.9.1) est un
morphisme propre de préschémas formels, comme il résulte de la définition et
de (II, 5.4.5).

Soient $\mathfrak X$ et $\mathfrak S$ deux préschémas formels localement
noethériens, $f:\mathfrak X\to\mathfrak S$ un morphisme \emph{de type fini}
(I, 10.13.3)~; les notations étant les mêmes que ci-dessus, on dit qu'une
partie $Z$ de l'espace sous-jacent à $\mathfrak X$ est \emph{propre sur}
$\mathfrak S$ (ou propre pour $f$) si, considérée comme partie de $X_0$, $Z$
est \emph{propre sur $S_0$} (II, 5.4.10). Toutes les propriétés des parties
propres de préschémas usuels énoncées dans (II, 5.4.10) sont encore valables
pour les parties propres de préschémas formels comme il résulte aussitôt des
définitions.
\end{env}

\begin{theorem}[3.4.2]
\label{III.3.4.2-fr}
Soient $\mathfrak X$, $\mathfrak Y$ deux préschémas formels localement
noethériens, $f:\mathfrak X\to\mathfrak Y$ un morphisme propre. Pour tout
$\mathcal O_{\mathfrak X}$-Module cohérent $\mathcal F$, les
$\mathcal O_{\mathfrak Y}$-Modules $R^qf_*(\mathcal F)$ sont cohérents pour
tout $q\geq0$.
\end{theorem}

Soient $\mathcal J$ un Idéal de définition de $\mathfrak Y$,
$\mathcal K=f^*(\mathcal J)\mathcal O_{\mathfrak X}$, et considérons les
$\mathcal O_{\mathfrak X}$-Modules
\[
\mathcal F_k=\mathcal F\otimes_{\mathcal O_{\mathfrak Y}}
(\mathcal O_{\mathfrak Y}/\mathcal J^{k+1})
=\mathcal F/\mathcal K^{k+1}\mathcal F\qquad(k\geq0)
\tag{3.4.2.1}\label{III.3.4.2.1-fr}
\]
qui forment évidemment un \emph{système projectif} de
$\mathcal O_{\mathfrak X}$-Modules topologiques, tel que

